% ==============================================================================
% Appendix: Agent-as-a-Judge Details
%
% Everything from the original appendix that is NOT in the main-paper section.
% The main paper (sec_judge_main_v2.tex) covers: motivation, human annotation /
% IAA, Agent-as-a-Judge pipeline description, and validation table.
% This appendix covers: design rationale, harness config, evidence package,
% rubric and iron rules, writing standard, automated checker, labeling
% protocol, regeneration loop, worked example, and full prompt template.
% ==============================================================================

\section{Agent-as-a-Judge Details}
\label{app:agent-as-a-judge}

This appendix documents the full implementation of the artifact-aware
Agent-as-a-Judge introduced in \S\ref{sec:judge}: an autonomous agent that
analyzes each trajectory for failure patterns across all \taxonomy lifecycle
stages and categorizes them according to the taxonomy. Each verdict is a
Markdown document, \texttt{analysis.md}, produced under a fixed rubric
covering the six lifecycle stages (A--F) plus cross-stage dynamics (X), and
passed through an automated structural/quantitative checker before being
accepted. Annotation proceeds in two phases. In the \emph{free-form} phase
the judge is given no pattern list and reports whatever defects it can
evidence; this is the phase from which \taxonomy was inductively derived
(\S\ref{sec:taxonomy}). In the \emph{labeling} phase each accepted issue
is mapped to an \taxonomy pattern ID and root-cause pillar
(\S\ref{app:evaluator:labeling}), and it is these labels that
Figure~\ref{fig:failure-attribution} aggregates. Below we document the
judge's execution harness, the evidence package it receives, the rubric
and scoring-integrity rules it is instructed to apply, the per-issue
writing standard, the automated checker, the labeling and human
calibration protocol, and the self-healing regeneration loop that retries
documents failing the checker.

\subsection{Design Rationale}
\label{app:evaluator:rationale}

Three properties motivate using an agentic judge rather than a static
rubric-matching classifier or a single LLM call on the transcript:

\begin{itemize}
  \item \textbf{Verification requires execution, not just reading.} Many
  failure modes in autonomous research rollouts (unit/magnitude bugs,
  degenerate optimization, evaluator-feedback fitting, silent train/test
  leakage, and computation run on synthetic or toy substrates) are only
  detectable by re-deriving a number from the rollout's own code and
  ground truth, not by pattern-matching the transcript text. The judge is
  therefore given the same execution environment as the rollout (a
  sandboxed shell) and is instructed to re-run cheap, dependency-light
  computations whenever a claim is both suspicious and not resolvable by
  inspection alone (\S\ref{app:evaluator:rubric}, Iron Rule~2).

  \item \textbf{Long, unstructured transcripts need a forcing function
  for coverage.} A single free-form ``critique this trajectory'' prompt
  reliably collapses onto whatever is most salient in the last few
  thousand tokens. We instead fix a stage skeleton
  (\S\ref{app:evaluator:rubric}) that the judge must fill in dedicated,
  quota-gated sections, using extraction tooling
  (\S\ref{app:evaluator:harness}) to navigate transcripts that can exceed
  4M characters without reading them in full.

  \item \textbf{A judge with no cross-checkable output is unfalsifiable.}
  Because the judge itself can hallucinate or under-verify, every document
  it produces is required to carry verifiable anchors (line numbers, file
  names, exact numeric values) for its claims, and is passed through an
  automated checker (\S\ref{app:evaluator:checker}) that specifically
  penalizes vague, anchor-free, or template-recycled writing before a
  document is accepted.
\end{itemize}

\subsection{Judge Execution Harness}
\label{app:evaluator:harness}

The judge is itself a Claude Code agent, invoked non-interactively with
full \texttt{Bash}/\texttt{Read}/\texttt{Write}/\texttt{Edit} access inside
a per-task scratch workspace, and \emph{no} network/search tools. Each
task is a fresh, zero-history session: the judge sees only what is placed
in its workspace for that one rollout, so its verdict cannot be
contaminated by patterns it inferred from other rollouts in the same
model/benchmark pool.

\begin{table}[h]
\centering
\small
\caption{Judge harness configuration. \bench comprises two pools
corresponding to the two task regimes: an
\emph{execution-feedback} pool (where a quantitative metric or human SOTA
exists) and a \emph{fully open-ended} pool (no external signal). Both share
this harness and differ only in pool-specific prompt injections}
\label{tab:harness-config}
\begin{tabular}{ll}
\toprule
Parameter & Value (execution-feedback pool) \\
\midrule
Judge model & claude-opus-5 (routed via OpenRouter) \\
Reasoning effort & \texttt{xhigh} \\
Max agent turns & 90 \\
Wall-clock budget per task & 5400\,s (90\,min) nominal;\\
 & extended to 12600\,s for tasks whose\\
 & verification re-runs need more headroom \\
Max regeneration attempts & 2 (per task, before the driver gives up;\\
 & overridable with \texttt{--force-tasks}) \\
Concurrency & 4--8 tasks in parallel \\
Directories mounted read-only & benchmark repository root; the source\\
 & run's task-definition / ground-truth root \\
\bottomrule
\end{tabular}

\end{table}

To make multi-megabyte execution logs tractable without truncating them,
the judge is given a small extraction CLI (\texttt{traj\_tools.py}) rather
than being told to parse logs itself:

\begin{verbatim}
traj_tools.py timeline    <log> --format <harness>   # tool-call timeline
traj_tools.py files       <log> --format <harness>   # every file written,
                                                       with provenance
traj_tools.py reconstruct <log> --format <harness> --name <fragment>
                                                      # replays Write+Edit
                                                      # history to reconstruct
                                                      # a file's final content
\end{verbatim}

The tool understands three distinct rollout-harness log formats (Claude
Code stream-JSON, Gemini CLI NDJSON, Codex CLI structured JSONL) and
normalizes them to a common timeline/file-provenance view. Reconstruction
from the log is treated as a way to narrate \emph{how} a file evolved
(false starts, abandoned versions); the judge is instructed that whenever
the rollout's real final container filesystem is available on disk, that
copy---not the log-reconstructed one---is definitive (Iron Rule~1,
\S\ref{app:evaluator:rubric}).

\subsection{Evidence Package Provided to the Judge}
\label{app:evaluator:evidence}

Artifact-aware annotation requires that the judge see what the rollout
\emph{made}, not only what it \emph{said}. For every (model, task) pair,
the driver therefore materializes a per-task judge workspace which contains evidence in Table~\ref{tab:evidence}.

\begin{table}[h]
\centering
\small
\caption{Evidence package materialized per task before the judge session
starts. The combination of execution log, delivered filesystem, and sealed
ground truth is what makes grounding failures such as C.1 (circular
validation on a synthetic substrate) detectable at all; none of them leaves
a signature in the final report.}
\label{tab:evidence}
\begin{tabular}{ll}
\toprule
Artifact & Role \\
\midrule
\texttt{problem\_readme.md}, & The task statement and data schema the\\
\texttt{data\_description.md} & rollout itself was given. \\
\texttt{<harness>.jsonl} & The rollout's complete execution log\\
 & (the primary evidence base). \\
\texttt{result.json} & Container-level execution status\\
 & (\texttt{status}/\texttt{duration}/\texttt{returncode}) only---\\
 & explicitly \emph{not} a quality signal. \\
\texttt{submissions.jsonl} & Every call the rollout made to the\\
 & benchmark's own scoring service, with\\
 & the score returned each time (when present). \\
\texttt{evaluator/} & The scoring service's real source code,\\
 & copied verbatim from the benchmark\\
 & definition---lets the judge check the\\
 & rollout's method against the exact metric\\
 & being computed, not just its stated name. \\
\texttt{verification.md} & Human-written notes (paper provenance,\\
 & held-out set construction, oracle score\\
 & ceiling) where available. \\
\texttt{agent\_code/} & The rollout's actual final container\\
 & filesystem (its real \texttt{workspace/}),\\
 & capped at 60 files / 1\,MB per file / 20\,MB\\
 & total, source-suffix filtered. \\
Sealed ground truth \& & Referenced by real host path (mounted\\
raw task data & read-only) rather than copied---these can\\
 & reach tens of GB per task---so the judge\\
 & can query them on demand without a full\\
 & read. \\
\bottomrule
\end{tabular}

\end{table}

\subsection{The Stage Rubric and Iron Rules}
\label{app:evaluator:rubric}

Every \texttt{analysis.md} must follow a fixed section skeleton (headings
reproduced verbatim; stage labels A--F and X correspond to the \taxonomy stage
axis of \S\ref{sec:taxonomy}):

\begin{quote}\small
\texttt{\# <title: task + one-line verdict>} \\
\texttt{> Core Verdict} (2--4 sentences, the 2--3 sharpest findings) \\
\texttt{\#\# Metadata} (task/model/harness, execution status, gold vs.\ agent observable) \\
\texttt{\#\# Trajectory Arc} (ideation $\to$ retrieval $\to$ execution $\to$ result $\to$ conclusion, one narrative paragraph) \\
\texttt{\#\# Credit Due} (genuine strengths---required, not optional) \\
\texttt{\#\# A. Ideation \& Planning} \\
\texttt{\#\# B. Retrieval \& Synthesis} \\
\texttt{\#\# C. Execution \& Implementation} (heaviest stage) \\
\texttt{\#\# D. Analysis \& Interpretation} \\
\texttt{\#\# E. Writing \& Documentation} \\
\texttt{\#\# F. Self-Verification \& Review} (guard against being fooled by an agent's own confident self-diagnosis) \\
\texttt{\#\# X. Cross-Stage Dynamics} (error propagation, goal drift, right-for-the-wrong-reason outcomes) \\
\texttt{\#\# Sentence-by-Sentence Checklist} (every key claim in the rollout's final report, marked pass / partial / fail) \\
\texttt{\#\# Numerical Grounding Notes} (what was independently re-derived, and the result) \\
\texttt{\#\# Retraction / Correction Log} (honest record of any self-corrected misjudgment) \\
\texttt{\#\# One-Line Verdict}
\end{quote}

\paragraph{Depth standard.} The judge is instructed to cover
\textbf{25--40 distinct issues} across all stages, each written as a
\emph{paragraph}, not a bullet: \textbf{mechanism} (what the rollout
concretely did) + \textbf{why it matters} + \textbf{the charitable/honest
reading} + \textbf{verifiable evidence} (a log line number, file name,
code identifier, or exact value). Each issue additionally carries a
one-line trailer
\texttt{[stage: <A--F,X> | root cause: <grounding | depth | integrity |
robustness>]}, giving the two coordinates of the \taxonomy axes. In the
free-form phase the judge assigns only these coordinates and never a
pattern ID, so that pattern boundaries are induced from the annotations
rather than presupposed by them. The calibration statistic is
characters-per-issue $\geq 280$ (target $\geq 290$) on the whole
document---a proxy for ``no issue is a one-line bullet.'' Genuinely limited
rollouts (a single fatal bug, an unrecoverable infrastructure failure, a
pure tautology) are explicitly allowed \emph{fewer} issues (16--20) as
long as each is written deep (characters-per-issue $\geq 350$): depth is
prioritized over hitting a raw count.

\paragraph{Scoring-signal literacy.} In the execution-feedback pool, where
the source benchmark provides a reward and a \texttt{reason} code, the
judge must not treat \texttt{reward=0} as a quality signal by default:

\begin{table}[h]
\centering
\small
\caption{Reward/reason taxonomy the judge must disambiguate before treating
a score as evidence of quality. This taxonomy is pool-specific; the
open-ended pool has no such field, and quality there is established purely
from internal evidence.}
\label{tab:rewardreason}
\begin{tabular}{lll}
\toprule
\texttt{reason} & Meaning & Quality signal? \\
\midrule
\texttt{judge\_unavailable} & Scoring infrastructure was down & No---rollout may be excellent \\
\texttt{no\_decision} & Malformed/incomplete decision artifact & No---a delivery failure, not a science failure \\
\texttt{soft[<observable>]} & Genuine score against the gold observable & Yes \\
\bottomrule
\end{tabular}

\end{table}

A related, high-frequency failure the judge is told to check explicitly
is \emph{observable mismatch}: the rollout computes a plausible, correctly
executed proxy quantity that is simply not the one the benchmark scores
against (e.g.\ computing group delay when the gold metric is efficiency),
which yields zero credit regardless of code quality.

\paragraph{Iron Rules (\textit{non-negotiable}).} These
codify lessons from earlier mis-judged trajectories in this project and
are stated to the judge verbatim:

\begin{enumerate}
  \item \textbf{Follow code evolution to the final delivered artifact.}
  Rollouts frequently contain multiple superseded versions (false starts,
  abandoned fallbacks). Never indict the final conclusion using a version
  the rollout itself discarded; when the final report explicitly states
  what it did, prefer that over a misleading earlier draft.

  \item \textbf{Sanity-check every task; re-run selectively, never
  exhaustively.} A near-zero-cost order-of-magnitude/units check catches
  most bugs (canonical examples we have caught this way: a protein--DNA
  interface burial reported as 47\,\AA$^2$ where $>1000$ is expected --
  almost certainly an nm$^2$/\AA$^2$ unit error; a water-dimer interaction
  energy of $-0.16$ where $\sim-5$ is expected). Full re-execution is
  reserved for cases where a number is both suspicious \emph{and}
  unresolvable by inspection \emph{and} cheap to reproduce
  (dependency-light, a few core lines, not the whole pipeline). Insight is
  not synonymous with re-running: many of the sharpest findings in this
  project were pure analytical derivations with zero re-execution.

  \item \textbf{Give credit where due.} Honestly reported null results,
  genuine mechanistic modeling, and self-caught bugs are real strengths
  that must be written up with the same rigor as failures.

  \item \textbf{Retract and log honestly.} If the judge itself misjudged
  something upon further reading, it must record the retraction and the
  lesson explicitly rather than silently editing it away.

  \item \textbf{Judge each trajectory on its own terms.} No
  cross-trajectory template conclusions.

  \item \textbf{Check for answer contamination.} If a rollout fetches the
  source paper's full text before modeling, its apparent ``independent
  replication'' of the paper's conclusion may simply be reading the answer
  first; any gold-adjacent fact that appears in text the rollout is shown
  to have already read cannot be credited as an independent finding
  (divergence from the paper, not mere disclosure, is what counts as
  evidence of independence).

  \item \textbf{Judge retrieval \emph{sufficiency}, not only retrieval
  \emph{honesty}.} A separate failure mode from fabrication/contamination
  is simply not searching for information the gold observable required
  (e.g.\ skipping a dataset search and using ``no data'' as an excuse to
  fall back to a synthetic model). This failure often disguises itself as
  a downstream ideation defect.

  \item \textbf{Verify every citation before crediting ``honest, no
  fabrication.''} Grep each cited DOI/arXiv ID/title against the actual
  retrieval log; an ID that never appears in a real search result but
  shows up in the final report is a hallucinated citation, not a
  formatting artifact.

  \item \textbf{Don't be disarmed by a rollout's own eloquent
  self-diagnosis.} Rollouts frequently \emph{name} their own core defect
  in a review paragraph and then ship the finding unchanged.
  ``Identified the mechanism'' $\neq$ ``addressed it.'' Credit for a
  review stage is recorded only against what the rollout
  \emph{independently found and then acted on}---naming a flaw without
  correcting it is a problem to flag, not a strength to credit. This rule
  is the operational definition behind pattern F.4.
\end{enumerate}

\subsection{Per-Issue Writing Standard}
\label{app:evaluator:writing}

Beyond the stage/quota structure, every numbered issue is required to
satisfy:

\begin{itemize}
  \item At least one \emph{verifiable anchor}---a concrete number, a log
  line index, a file name, or a code identifier---per issue.
  \item A well-formed \texttt{[stage | root cause]} trailer, since
  downstream aggregation into
  Figure~\ref{fig:failure-attribution} depends on it.
  \item A minimum length ($\geq 200$ characters); single-sentence bullets
  are rejected or must be merged into a fuller issue.
  \item No large verbatim pasting from \texttt{problem\_readme.md} or
  \texttt{result.json}---pasting is explicitly not analysis and is
  penalized as padding.
  \item No templated/recycled phrasing across issues---every issue must
  state a fact unique to that specific rollout.
\end{itemize}

\subsection{Automated Checker Script (\texttt{qa\_check\_analysis.py})}
\label{app:evaluator:checker}

Every generated \texttt{analysis.md} is passed through a deterministic
Python checker before being accepted, rather than relying on a second LLM
pass to judge the first judge. The checker enforces six independent gate
classes, calibrated against the p10--p25 percentiles of a 134-document
reference corpus (deliberately below the median, so as not to reject the
same style of document the rubric targets):

\begin{enumerate}
  \item \textbf{Totals.} Whole-document character count, total numbered
  issues, and characters-per-issue, each with a floor (nominal:
  $\geq 16{,}000$ chars, $\geq 28$ issues, $\geq 300$ chars/issue; floors
  scale to $0.6\times$ for pools/reasons that are naturally thin, e.g.\
  infrastructure-failure trajectories).

  \item \textbf{Section length.} A minimum character floor per required
  section (Table~\ref{tab:sectionfloors}), so no single stage can be a
  one-line stub while another stage is padded.

  \item \textbf{Section breadth.} A minimum issue count for each stage
  section (A: $\geq 3$, B: $\geq 4$, C: $\geq 10$, D: $\geq 3$, E:
  $\geq 3$, F: $\geq 3$, X: $\geq 2$), so genuine coverage cannot be
  concentrated in one stage.

  \item \textbf{Issue thickness distribution.} The \emph{distribution} of
  per-issue length is checked, not just its mean: documents are rejected
  if more than 28\% of issues fall under 200 characters or more than 5\%
  fall under 100 characters.

  \item \textbf{Density / anti-padding.} A minimum count of
  \emph{verifiable anchors} (code spans, numeric literals, file/line
  references, DOI/arXiv identifiers) across the document, a minimum
  fraction of issues carrying at least one anchor ($\geq 85\%$), a cap on
  near-duplicate 60-character text shingles (catches templated filler,
  $\leq 8\%$), and---when the task workspace is available to the
  checker---a cap on verbatim-copied spans traced back to the source
  artifacts ($\leq 20\%$ of document length) alongside a minimum count of
  quoted spans that \emph{do} verifiably trace back to a real source file
  (rewarding grounded quotation while penalizing wholesale copying).

  \item \textbf{Label well-formedness.} Every numbered issue must carry a
  parseable \texttt{[stage | root cause]} trailer with values drawn from
  the fixed vocabularies; documents with unparseable or out-of-vocabulary
  trailers are rejected, since stage-level attribution depends on them.
\end{enumerate}

\begin{table}[h]
\centering
\small
\caption{Per-section character floors enforced by the automated checker.}
\label{tab:sectionfloors}
\begin{tabular}{lr}
\toprule
Section & Min.\ characters \\
\midrule
Core Verdict & 850 \\
Metadata & 720 \\
Trajectory Arc & 380 \\
Credit Due & 760 \\
A. Ideation \& Planning & 760 \\
B. Retrieval \& Synthesis & 1020 \\
C. Execution \& Implementation & 2550 \\
D. Analysis \& Interpretation & 600 \\
E. Writing \& Documentation & 640 \\
F. Self-Verification \& Review & 640 \\
X. Cross-Stage Dynamics & 520 \\
Sentence-by-Sentence Checklist (min.\ 12 rows) & 1100 \\
Numerical Grounding Notes & 600 \\
Retraction / Correction Log & 340 \\
One-Line Verdict & 210 \\
\bottomrule
\end{tabular}

\end{table}

A document that fails any gate is reported with the exact list of failed
gates (e.g.\ \texttt{section too short:Analysis 480 < 600},
\texttt{section too few issues:Execution 7 < 10}); this failure report is
what drives the regeneration loop below.

\subsection{\taxonomy Labeling and Human Calibration}
\label{app:evaluator:labeling}

Accepted documents are converted to structured labels in a second pass.
Each numbered issue, together with its evidence anchor and its
\texttt{[stage | root cause]} trailer, is presented to a labeling judge
holding the \taxonomy pattern definitions, which assigns exactly one pattern ID
or \texttt{none} where no pattern fits. The \texttt{none} rate was tracked throughout taxonomy development:
patterns were added whenever a recurring unlabeled behavior emerged, and the
taxonomy was frozen once further annotation produced no new pattern candidates
(theoretical saturation, \S\ref{sec:judge:human}).

A trajectory may exhibit multiple patterns, and each cell of
Figure~\ref{fig:failure-attribution} counts distinct trajectories rather
than issues.

Calibration proceeds against human annotation on a stratified sample of
50 trajectories spanning all eight model--harness combinations. Three
experts independently annotated each sampled trajectory under the same
rubric (inter-annotator agreement $\kappa = 0.85$); judge--human
agreement at both the pattern and taxonomy-categorization levels is
reported in \S\ref{sec:judge:validation} and
Table~\ref{tab:judge-performance}. These figures are aggregate; we do not
report per-pillar agreement at this time. Patterns anchored to concrete artifacts
(R1, R4) are more reliably adjudicated than patterns requiring a
judgment about the agent's reasoning (R2), so we treat R2 rates as
lower-confidence throughout the analysis.

\subsection{Self-Healing Regeneration Loop}
\label{app:evaluator:loop}

The driver (\texttt{generate\_analysis\_cc.py}) re-globs the benchmark's
result directory on every pass, so newly completed rollouts are picked up
automatically. For each (model, task):

\begin{enumerate}
  \item Skip if an accepted \texttt{analysis.md} already exists
  (\texttt{--resume}).

  \item Otherwise launch a fresh judge session
  (\S\ref{app:evaluator:harness}); on completion, run the checker
  (\S\ref{app:evaluator:checker}).

  \item If the checker fails, the exact failed-gate list is injected back
  into the next attempt's prompt with gate-specific remediation guidance
  (e.g.\ \texttt{section too few issues} $\to$ ``go find real additional
  problems in that stage, don't split one existing issue into three'';
  \texttt{near-duplicate share} $\to$ ``delete the repeated boilerplate,
  write specifics''), together with a hard requirement that the new
  attempt have \emph{strictly more} issues than the last failed attempt
  ($\geq$ previous $+5$) without discarding any previously-correct
  finding.

  \item After a fixed number of failed attempts on the same task, the
  driver gives up on that task for the current pass (a cost
  circuit-breaker) rather than silently retrying forever; a task can be
  forced past this cap for a manual follow-up pass.

  \item The loop repeats across the whole pool until every discoverable
  task has an accepted document or the pass budget is exhausted.
\end{enumerate}

In practice, most failed first attempts fail on one of two patterns: (a)
a borderline miss on a single length/issue-count gate, resolved by a
fresh attempt; or (b) the judge's own in-session verification re-run
(Iron Rule~2) exceeding the wall-clock budget before it can write up the
remaining sections, resolved by extending the per-task timeout rather
than by discarding the verification step.

\subsection{Worked Example}
\label{app:evaluator:example}

The following (translated, lightly condensed) excerpt from an accepted
\texttt{analysis.md} (task \texttt{echonet\_lvef}, model \texttt{sonnet-5})
illustrates the mechanism+harm+charitable-reading+evidence issue format
required by \S\ref{app:evaluator:writing}. The rollout designed a correct
modeling approach, launched training in the background, and then ended its
turn relying on an unverified scheduled-wakeup promise; the container was
torn down before training finished, so zero predictions were ever
submitted. The issue was labeled \texttt{[stage: X | root cause:
robustness]} in the free-form phase and mapped to X.8 in the labeling pass.

\begin{quote}\small
\textbf{Core Verdict (excerpt).} ``This is a trajectory with \emph{zero
delivery}, not a trajectory with poor scientific quality, and the two
must be accounted for separately. \texttt{result.json}'s
\texttt{status=success}, \texttt{returncode=0} only mean the container
process exited cleanly\ldots\ the final \texttt{workspace/} contains only
five files\ldots\ and a 0-byte \texttt{train.log}. [\ldots] The agent
trusted the promise returned by \texttt{ScheduleWakeup}---`the harness
will re-invoke you'---but the session ended before that wakeup ever fired,
and the container, along with the training process, was destroyed with it.
It outsourced its delivery responsibility to a tool guarantee it never
verified and which was not even in its allow-listed tool set.''
\end{quote}

\begin{quote}\small
\textbf{Credit Due (excerpt).} ``Its engineering validation was not
perfunctory---it actually measured what it claimed. Three separate
timing passes make the point: a synthetic-tensor benchmark first measured
14.67\,s/step; rather than concluding `too slow, switch models,' it
correctly identified this as CUDA kernel-compile warmup, re-measured at
0.1007\,s/step after warmup, and then \emph{did not stop at that
favorable synthetic number}---it re-measured end-to-end on real data
through the full DataLoader and augmentation path, got 0.2156\,s/batch,
and revised its own throughput estimate downward by $2.1\times$
accordingly. This is the most researcher-like moment in the trajectory:
distinguishing microbenchmark from end-to-end throughput, and defaulting
to the more conservative number.''
\end{quote}

\subsection{Full Judge Prompt Template}
\label{app:evaluator:fulltemplate}

For exact reproducibility we reproduce the structure of the
per-task instruction below.

\begin{lstlisting}[style=prompt]
# Task: write a comprehensive, information-dense analysis.md for ONE
# agentic scientific-ML trajectory.

You are doing a deep trajectory audit of one agentic-scientific-discovery
rollout. This session analyzes exactly ONE trajectory -- judge it on its
own terms, do not apply cross-trajectory template conclusions. You are
fully autonomous: do not ask questions, work through to a finished
deliverable.

[If this is a redo after a failed quality gate: the exact list of failed
 gates from the previous attempt is inserted here, with the hard
 requirement that the new issue count be >= previous + 5.]

## 0. Required reading (read fully before doing anything else)
1. Framework: {onboarding path} -- follow its workflow, depth standard,
   six-stage structure, and Iron Rules exactly. [Pool-specific
   suppressions/overrides of individual Iron Rules are noted here --
   e.g. for a pool with no retrieval tools, the citation/contamination
   Iron Rules are marked not applicable.]
2. [Optional] a worked exemplar analysis.md the required depth should match.

## 1. This task's data (all in the current working directory)
- meta.json -- task_id, model, harness format.
- problem_readme.md -- the task statement the rollout itself received.
- data_description.md -- data schema, if present.
- result.json -- container execution status ONLY (status/duration/
  returncode). This is NOT a quality signal -- do not conflate it with
  a reward/reason field from a different pool.
- [pool-specific scoring-integrity block -- see below]
- evaluator/ -- the scoring service's real source, verbatim. Read it
  line by line: figure out exactly what it scores, at what tolerance,
  before judging whether the rollout's method actually matches that
  logic or merely guessed the right output format.
- verification.md -- human-written verification notes, if present (paper
  provenance, held-out construction, oracle score ceiling). Read this
  first; it often names the single most common systematic error on this
  task.
- sealed ground truth -- referenced at {real host path} via a read-only
  mount, NOT copied (can be 100s of MB to tens of GB). Query it
  selectively; never read it in full.
- agent_code/ -- the container's actual final workspace, {N} files,
  {M} skipped as non-source/oversized (real path: {src}). This is the
  authoritative delivered code; log reconstruction is only used to
  narrate how the rollout arrived at it.
- {harness}.jsonl -- the rollout's full execution log ({K}k characters).
  DO NOT read it in full -- use traj_tools.py timeline/files/reconstruct
  to navigate, then sed -n/grep -n for the specific spans you need.

## 2. [Pool-specific integrity section]
[\bench pool:] There is no reward/reason taxonomy in this pool --
quality judgment rests entirely on (a) internal evidence: code,
execution log, numerical self-consistency, evaluator source logic; and
(b) submissions.jsonl (if present) -- the real scored-submission history.
Do not invent or borrow a reason field from a different pool's schema.

## 3. Iron Rules for this task
[Pool-specific selection/adaptation of the shared Iron Rules from
 ONBOARDING(2).md -- e.g. for the \bench pool: retrieval-integrity
 rules are marked not applicable (no retrieval tools exist in this
 container); the evaluator-feedback-fitting rule (see
 Sec.~\ref{app:evaluator:adaptation}) is elevated to the single most
 important check for this task and MUST be written up as its own
 standalone, evidenced issue.]
- Follow multi-version code to the final delivered artifact.
- Don't be fooled by a beautiful self-diagnosis that was never acted on.
- Perfect/extreme scores get checked too: decompose them; check whether
  they are propped up by privileged information (e.g. an evaluator
  default value, a leaked starting condition) rather than a correct method.

## 4. Output requirements -- hard gates, self-check before you finish
### 4.1 Skeleton (all sections required, headings copied verbatim)
[the fourteen headings from Sec.~\ref{app:evaluator:rubric}]
### 4.2 Quota table (targets -- calibrated above the automated floor)
[per-section target character count and target issue count --
 see Table~\ref{tab:sectionfloors} for the corresponding pass/fail floor]
### 4.3 How to write each issue (this is what "information density" means)
Each issue = mechanism (what it concretely did) + why it's harmful +
the honest/charitable reading + evidence, written as one full paragraph.
- At least one verifiable anchor per issue: a number, a log line index,
  a file name, a code identifier.
- A one-sentence bullet does not count as an issue. Merge or expand
  anything under ~200 characters.
- Do not paste large verbatim spans from problem_readme.md/result.json --
  pasting is not analysis and will be scored as padding.
- Do not recycle templated phrasing -- every issue must state a fact
  unique to this specific trajectory.
### 4.4 Give full credit
Genuine strengths must be written up with the same rigor as failures:
honestly reported below-target results, genuinely independent
verification design (held-out/cross-validation), self-caught bugs,
honestly-flagged failed submissions, reproducible seeds and numbers.

## 5. Where to write it
{out_dir}/analysis.md (directory already created; write ONLY this file).
State the one-line verdict in the H1 title itself.
Before finishing, self-check: word count, issue count, that none of the
six stages is perfunctory, that the checklist has enough rows. If it does
not meet the bar, keep digging -- do not submit early.

## 6. Boundaries
Deliver only this one file. Do not modify anything else in the workspace,
do not refactor anything. Do all verification/re-computation in your own
main loop -- do not spawn sub-agents. You may independently re-run
scoring logic locally when its dependencies are light (numpy/scipy/sympy-
level); ground-truth/raw-data directories are read-only and must be
queried selectively, never traversed in full.
\end{lstlisting}